%%%%%%%%%%%%%%%%%%%%%%%%%%%%%%%%%%%%%%%%%%%%%%%%%%%%%%%%%%%%%%%%%%%
% Standalone general university LaTeX template -- homework sheet
% One-file version: no \input files needed.
%%%%%%%%%%%%%%%%%%%%%%%%%%%%%%%%%%%%%%%%%%%%%%%%%%%%%%%%%%%%%%%%%%%

\documentclass[11pt,a4paper]{scrartcl}

%%%%%%%%%%%%%%%%%%%%%%%%%%%%%%%%%%%%%%%%%%%%%%%%%%%%%%%%%%%%%%%%%%%
% Embedded file: includes/university_metadata.tex
%%%%%%%%%%%%%%%%%%%%%%%%%%%%%%%%%%%%%%%%%%%%%%%%%%%%%%%%%%%%%%%%%%%
\newcommand{\CourseTitle}{COURSE NAME}
\newcommand{\CourseShortTitle}{COURSE SHORT NAME}
\newcommand{\DocumentKind}{DOCUMENT KIND}
\newcommand{\DocumentTitle}{DOCUMENT TITLE}
\newcommand{\DocumentSubtitle}{}
\newcommand{\AuthorName}{YOUR NAME}
\newcommand{\InstitutionName}{INSTITUTION NAME}
\newcommand{\SubmissionDate}{\today}

\newcommand{\makeuniversitytitle}{%
    \begin{titlepage}
        \centering
        \vspace*{2cm}
        {\Large \CourseTitle\par}
        \vspace{1.2cm}
        {\huge\bfseries \DocumentTitle\par}
        \ifx\DocumentSubtitle\empty\else
            \vspace{0.4cm}
            {\Large \DocumentSubtitle\par}
        \fi
        \vspace{1cm}
        {\Large \DocumentKind\par}
        \vfill
        {\large \AuthorName\par}
        \vspace{0.2cm}
        {\large \InstitutionName\par}
        \vspace{0.8cm}
        {\large \SubmissionDate\par}
    \end{titlepage}
    \pagenumbering{arabic}
}

\newcommand{\makechaptertitle}{%
    \begin{center}
        {\Large \CourseTitle\par}
        \vspace{0.5cm}
        {\huge\bfseries \DocumentTitle\par}
        \ifx\DocumentSubtitle\empty\else
            \vspace{0.25cm}
            {\Large \DocumentSubtitle\par}
        \fi
        \vspace{0.5cm}
        {\large \AuthorName \hfill \SubmissionDate\par}
    \end{center}
    \vspace{0.5cm}
}

\newcommand{\makephysiktitle}{\makeuniversitytitle}

%%%%%%%%%%%%%%%%%%%%%%%%%%%%%%%%%%%%%%%%%%%%%%%%%%%%%%%%%%%%%%%%%%%
% Document metadata
%%%%%%%%%%%%%%%%%%%%%%%%%%%%%%%%%%%%%%%%%%%%%%%%%%%%%%%%%%%%%%%%%%%
\renewcommand{\CourseTitle}{Physik IV: Kerne und Teilchen}
\renewcommand{\CourseShortTitle}{Physik IV}
\renewcommand{\DocumentKind}{Hausaufgabe / Tafelvortrag / Lernfassung}
\renewcommand{\DocumentTitle}{Aufgabenblatt 10}
\renewcommand{\DocumentSubtitle}{Elektromagnetische Schauer, Cherenkov-Licht und Bethe-Bloch-Formel}
\renewcommand{\AuthorName}{Tobias Laser}
\renewcommand{\InstitutionName}{Universit\"at Innsbruck}
\renewcommand{\SubmissionDate}{30. Mai 2026}

%%%%%%%%%%%%%%%%%%%%%%%%%%%%%%%%%%%%%%%%%%%%%%%%%%%%%%%%%%%%%%%%%%%
% Embedded file: includes/university_packages.tex
%%%%%%%%%%%%%%%%%%%%%%%%%%%%%%%%%%%%%%%%%%%%%%%%%%%%%%%%%%%%%%%%%%%
\usepackage[margin=2.5cm]{geometry}
\usepackage[onehalfspacing]{setspace}
\usepackage{scrlayer-scrpage}

\usepackage[T1]{fontenc}
\usepackage[utf8]{inputenc}
\usepackage[ngerman]{babel}
\usepackage{lmodern}
\usepackage{microtype}
\usepackage{csquotes}

\usepackage{amsmath,amsthm,amssymb,mathtools}
\usepackage{bm}
\usepackage{icomma}
\usepackage{cancel}

\usepackage{graphicx}
\usepackage{tikz}
\usetikzlibrary{arrows.meta,calc,decorations.pathmorphing,positioning,patterns,angles,quotes}
\usepackage{pgfplots}
\pgfplotsset{compat=1.18}

\usepackage[locale=DE,space-before-unit=true,per-mode=symbol,separate-uncertainty=true]{siunitx}
\usepackage{booktabs,multirow,array,tabularx,longtable}
\usepackage{caption,subcaption}
\usepackage{placeins}
\usepackage{enumitem}
\usepackage{xparse}

\usepackage[most]{tcolorbox}
\tcbuselibrary{breakable,skins,theorems}

\usepackage[breaklinks=true,colorlinks=true,linkcolor=blue,urlcolor=blue,citecolor=blue]{hyperref}
\usepackage[nameinlink,noabbrev]{cleveref}

\clearpairofpagestyles
\ihead{\CourseShortTitle}
\ohead{\DocumentKind}
\cfoot{\pagemark}
\pagestyle{scrheadings}

\setlist[enumerate]{itemsep=0.4em,topsep=0.4em}
\setlist[itemize]{itemsep=0.25em,topsep=0.35em}
\captionsetup{font=small,labelfont=bf}

%%%%%%%%%%%%%%%%%%%%%%%%%%%%%%%%%%%%%%%%%%%%%%%%%%%%%%%%%%%%%%%%%%%
% Embedded file: includes/university_macros.tex
%%%%%%%%%%%%%%%%%%%%%%%%%%%%%%%%%%%%%%%%%%%%%%%%%%%%%%%%%%%%%%%%%%%
\newcommand{\R}{\mathbb{R}}
\newcommand{\N}{\mathbb{N}}
\newcommand{\Z}{\mathbb{Z}}
\newcommand{\Q}{\mathbb{Q}}
\newcommand{\C}{\mathbb{C}}

\newcommand{\set}[1]{\left\{ #1 \right\}}
\newcommand{\abs}[1]{\left| #1 \right|}
\newcommand{\norm}[1]{\left\lVert #1 \right\rVert}
\newcommand{\avg}[1]{\left\langle #1 \right\rangle}
\newcommand{\paren}[1]{\left( #1 \right)}
\newcommand{\bracket}[1]{\left[ #1 \right]}
\newcommand{\ol}[1]{\overline{#1}}

\newcommand{\dd}{\,\mathrm{d}}
\newcommand{\dv}[2]{\frac{\dd #1}{\dd #2}}
\newcommand{\pdv}[2]{\frac{\partial #1}{\partial #2}}
\newcommand{\pddv}[2]{\frac{\partial^2 #1}{\partial #2^2}}
\newcommand{\evalat}[2]{\left. #1 \right|_{#2}}

\newcommand{\vect}[1]{\bm{#1}}
\newcommand{\unitvect}[1]{\hat{\bm{#1}}}
\newcommand{\grad}{\vect{\nabla}}
\newcommand{\divergence}{\grad\!\cdot}
\newcommand{\curl}{\grad\!\times}

\newcommand{\half}{\frac{1}{2}}
\newcommand{\third}{\frac{1}{3}}
\newcommand{\twothirds}{\frac{2}{3}}
\newcommand{\hbarc}{\hbar c}
\newcommand{\alphaf}{\alpha}
\newcommand{\eV}{\si{eV}}
\newcommand{\MeV}{\si{MeV}}
\newcommand{\GeV}{\si{GeV}}
\newcommand{\TeV}{\si{TeV}}

\newcommand{\nuclide}[3]{^{#1}_{#2}\mathrm{#3}}
\newcommand{\isotope}[2]{^{#1}\mathrm{#2}}
\newcommand{\anti}[1]{\overline{#1}}
\newcommand{\pbar}{\anti{p}}
\newcommand{\nbar}{\anti{n}}
\newcommand{\ee}{e^-}
\newcommand{\ep}{e^+}
\newcommand{\mup}{\mu^+}
\newcommand{\mum}{\mu^-}
\newcommand{\nue}{\nu_e}
\newcommand{\numu}{\nu_\mu}
\newcommand{\nutau}{\nu_\tau}
\newcommand{\pionp}{\pi^+}
\newcommand{\pionm}{\pi^-}
\newcommand{\pionz}{\pi^0}
\newcommand{\kaonp}{K^+}
\newcommand{\kaonm}{K^-}
\newcommand{\kaonz}{K^0}

\newcommand{\diffxs}{\frac{\dd\sigma}{\dd\Omega}}
\newcommand{\totalxs}{\sigma_{\mathrm{total}}}
\newcommand{\lum}{\mathcal{L}}
\newcommand{\smand}{s}
\newcommand{\tmand}{t}
\newcommand{\umand}{u}
\newcommand{\Ecm}{E_{\mathrm{CM}}}

\newcommand{\ie}{d.\,h.}
\newcommand{\eg}{z.\,B.}
\newcommand{\wrt}{bez\"uglich}

\DeclareSIUnit{\barn}{b}
\DeclareSIUnit{\millibarn}{mb}
\DeclareSIUnit{\microbarn}{\micro b}
\DeclareSIUnit{\nanobarn}{nb}
\DeclareSIUnit{\fermi}{fm}
\DeclareSIUnit{\year}{a}

%%%%%%%%%%%%%%%%%%%%%%%%%%%%%%%%%%%%%%%%%%%%%%%%%%%%%%%%%%%%%%%%%%%
% Embedded file: includes/university_boxes.tex
%%%%%%%%%%%%%%%%%%%%%%%%%%%%%%%%%%%%%%%%%%%%%%%%%%%%%%%%%%%%%%%%%%%
\tcbset{
    unibox/.style={
        enhanced,
        breakable,
        sharp corners,
        boxrule=0.6pt,
        left=1.2mm,
        right=1.2mm,
        top=1mm,
        bottom=1mm,
        before skip=0.8em,
        after skip=0.8em,
        fonttitle=\bfseries
    }
}

\newtcolorbox{calculationbox}[1][]{unibox,title=Rechnung,colback=black!2,colframe=black!55,#1}
\newtcolorbox{sidecalcbox}[1][]{unibox,title=Nebenrechnung,colback=black!1,colframe=black!40,#1}
\newtcolorbox{solutionbox}[1][]{unibox,title=L\"osung,colback=blue!2,colframe=blue!45!black,#1}
\newtcolorbox{answerbox}[1][]{unibox,title=Antwort,colback=green!3,colframe=green!45!black,#1}
\newtcolorbox{notebox}[1][]{unibox,title=Notiz,colback=yellow!6,colframe=yellow!45!black,#1}
\newtcolorbox{definitionbox}[1][]{unibox,title=Definition,colback=cyan!3,colframe=cyan!45!black,#1}
\newtcolorbox{warningbox}[1][]{unibox,title=Achtung,colback=red!3,colframe=red!55!black,#1}
\newtcolorbox{summarybox}[1][]{unibox,title=Zusammenfassung,colback=violet!3,colframe=violet!50!black,#1}
\newtcolorbox{formulabox}[1][]{unibox,title=Formel,colback=orange!4,colframe=orange!65!black,#1}
\newtcolorbox{taskbox}[1][]{unibox,title=Aufgabe,colback=white,colframe=black!75,#1}

%%%%%%%%%%%%%%%%%%%%%%%%%%%%%%%%%%%%%%%%%%%%%%%%%%%%%%%%%%%%%%%%%%%
% Embedded file: includes/university_theorems.tex
%%%%%%%%%%%%%%%%%%%%%%%%%%%%%%%%%%%%%%%%%%%%%%%%%%%%%%%%%%%%%%%%%%%
\theoremstyle{definition}
\newtheorem{definition}{Definition}[section]
\newtheorem{example}{Beispiel}[section]
\theoremstyle{plain}
\newtheorem{satz}{Satz}[section]
\newtheorem{lemma}{Lemma}[section]
\theoremstyle{remark}
\newtheorem*{remark}{Bemerkung}

%%%%%%%%%%%%%%%%%%%%%%%%%%%%%%%%%%%%%%%%%%%%%%%%%%%%%%%%%%%%%%%%%%%
\begin{document}

\makeuniversitytitle
\tableofcontents
\newpage

%%%%%%%%%%%%%%%%%%%%%%%%%%%%%%%%%%%%%%%%%%%%%%%%%%%%%%%%%%%%%%%%%%%
\section{Aufgabe 1: Elektromagnetische Teilchenschauer im Heitler-Modell}

\begin{taskbox}
Ein hochenergetisches Elektron oder Photon l\"ost in der Atmosph\"are einen elektromagnetischen Schauer aus. Im vereinfachten Heitler-Modell wird nach jeder Strahlungsl\"ange $X_0$ die Teilchenzahl verdoppelt und die Energie gleichm\"a\ss{}ig auf die Teilchen verteilt. Der Schauer entwickelt sich, bis die Energie pro Teilchen die kritische Energie $E_c$ erreicht.
\end{taskbox}

\begin{definitionbox}[title=Physikalische Idee]
Ein elektromagnetischer Schauer wird im Wesentlichen durch zwei Prozesse getragen:
\begin{itemize}
    \item \textbf{Bremsstrahlung:} $e^{\pm}$ strahlen im Coulombfeld der Luftmolek\"ule Photonen ab.
    \item \textbf{Paarbildung:} hochenergetische Photonen erzeugen $e^-e^+$-Paare.
\end{itemize}
Das Heitler-Modell ersetzt die realistische stochastische Entwicklung durch ein deterministisches Verdopplungsmodell.
\end{definitionbox}

\subsection*{(a) Erste drei Stufen des Schauers}

\begin{calculationbox}[title=Skizze als Baumdiagramm]
Die durchgezogenen Linien stehen f\"ur Leptonen $e^-$ oder $e^+$; gewellte Linien stehen f\"ur Photonen $\gamma$. Die horizontale Achse ist hier die durchlaufene atmosph\"arische Tiefe in Einheiten von $X_0$.

\begin{center}
\begin{tikzpicture}[x=3.0cm,y=0.8cm,>=Latex]
    \foreach \x/\lab in {0/$0$,1/$X_0$,2/$2X_0$,3/$3X_0$}{
        \draw[gray!45] (\x,-3.1)--(\x,3.1);
        \node[above] at (\x,3.1) {\lab};
    }
    \node[left] at (0,0) {$e^-(E_0)$};

    \draw[thick,->] (0,0)--(1,0.9) node[midway,above] {$e^-$};
    \draw[decorate,decoration={snake,amplitude=1.3pt,segment length=5pt},->] (0,0)--(1,-0.9) node[midway,below] {$\gamma$};

    \draw[thick,->] (1,0.9)--(2,1.9) node[midway,above] {$e^-$};
    \draw[decorate,decoration={snake,amplitude=1.3pt,segment length=5pt},->] (1,0.9)--(2,0.9) node[midway,above] {$\gamma$};
    \draw[thick,->] (1,-0.9)--(2,-0.35) node[midway,above] {$e^-$};
    \draw[thick,->] (1,-0.9)--(2,-1.55) node[midway,below] {$e^+$};

    \draw[thick,->] (2,1.9)--(3,2.6) node[midway,above] {$e^-$};
    \draw[decorate,decoration={snake,amplitude=1.3pt,segment length=5pt},->] (2,1.9)--(3,1.8) node[midway,below] {$\gamma$};
    \draw[thick,->] (2,0.9)--(3,1.1) node[midway,above] {$e^-$};
    \draw[thick,->] (2,0.9)--(3,0.3) node[midway,below] {$e^+$};
    \draw[thick,->] (2,-0.35)--(3,-0.35) node[midway,above] {$e^-$};
    \draw[decorate,decoration={snake,amplitude=1.3pt,segment length=5pt},->] (2,-0.35)--(3,-1.0) node[midway,below] {$\gamma$};
    \draw[thick,->] (2,-1.55)--(3,-1.8) node[midway,below] {$e^+$};
    \draw[decorate,decoration={snake,amplitude=1.3pt,segment length=5pt},->] (2,-1.55)--(3,-2.6) node[midway,below] {$\gamma$};
\end{tikzpicture}
\end{center}
\end{calculationbox}

\begin{notebox}
Diese Skizze ist kein Feynman-Diagramm. Sie ist ein \textbf{Entwicklungsbaum} des Heitler-Modells. Die relevante Z\"ahlregel lautet: Nach $n$ Schritten gibt es idealisiert $2^n$ Teilchen mit ungef\"ahr gleicher Energie.
\end{notebox}

\subsection*{(b) Schauermaximum $X_{\max}$ und maximale Teilchenzahl $N_{\max}$}

\begin{calculationbox}
Nach $n$ Schritten gilt im Heitler-Modell
\[
    N(n)=2^n,
    \qquad
    E(n)=\frac{E_0}{2^n}.
\]
Das Schauermaximum ist erreicht, wenn die Energie pro Teilchen etwa die kritische Energie erreicht:
\[
    E(n_{\max})=E_c.
\]
Damit folgt
\[
    \frac{E_0}{2^{n_{\max}}}=E_c
    \quad\Longrightarrow\quad
    2^{n_{\max}}=\frac{E_0}{E_c}
    \quad\Longrightarrow\quad
    n_{\max}=\frac{\ln(E_0/E_c)}{\ln 2}.
\]
Da jeder Schritt einer Strahlungsl\"ange $X_0$ entspricht, ist
\[
    X_{\max}=n_{\max}X_0.
\]
\end{calculationbox}

\begin{answerbox}
\[
    \boxed{\frac{X_{\max}}{X_0}=\frac{\ln(E_0/E_c)}{\ln 2}},
    \qquad
    \boxed{N_{\max}=2^{n_{\max}}=\frac{E_0}{E_c}}.
\]
\end{answerbox}

\subsection*{(c) Schauerl\"angen f\"ur $E_0=10^{11}\,\si{eV}$ und $E_0=10^{16}\,\si{eV}$}

\begin{calculationbox}
Die kritische Energie ist
\[
    E_c=\SI{84}{MeV}=8.4\cdot 10^7\,\si{eV}.
\]
Das Schauermaximum liegt laut Angabe ungef\"ahr bei zwei Dritteln der gesamten Schauerl\"ange $L$:
\[
    X_{\max}\approx \frac{2}{3}L
    \quad\Longrightarrow\quad
    L\approx \frac{3}{2}X_{\max}.
\]
\end{calculationbox}

\begin{calculationbox}[title=Numerik]
F\"ur $E_0=10^{11}\,\si{eV}$ gilt
\[
    \frac{X_{\max}}{X_0}
    =\frac{\ln\paren{10^{11}/(8.4\cdot 10^7)}}{\ln 2}
    =10.22,
\]
also
\[
    \frac{L}{X_0}\approx \frac{3}{2}\cdot 10.22=15.33.
\]
F\"ur $E_0=10^{16}\,\si{eV}$ gilt
\[
    \frac{X_{\max}}{X_0}
    =\frac{\ln\paren{10^{16}/(8.4\cdot 10^7)}}{\ln 2}
    =26.83,
\]
also
\[
    \frac{L}{X_0}\approx \frac{3}{2}\cdot 26.83=40.24.
\]
\end{calculationbox}

\begin{answerbox}
\[
\begin{array}{c|c|c}
E_0 & X_{\max}/X_0 & L/X_0 \\
\hline
10^{11}\,\si{eV} & 10.22 & 15.33 \\
10^{16}\,\si{eV} & 26.83 & 40.24
\end{array}
\]
Die kritische Energie wird dadurch bestimmt, dass die dominierenden Verlustmechanismen wechseln: Oberhalb von $E_c$ dominieren Bremsstrahlung und Paarbildung; unterhalb von $E_c$ werden Ionisation und Anregung des Mediums wichtiger, sodass die weitere Teilchenproduktion ineffizient wird.
\end{answerbox}

\subsection*{(d) H\"ohe des Schauermaximums}

\begin{calculationbox}
Die atmosph\"arische Tiefe wird beschrieben durch
\[
    X=X_E e^{-h/H_0},
    \qquad
    H_0=\SI{8}{km},
    \qquad
    X_E=\SI{1030}{g/cm^2}.
\]
Aufgel\"ost nach $h$:
\[
    h=-H_0\ln\paren{\frac{X}{X_E}}.
\]
Die Tiefe des Maximums ist
\[
    X_{\max}^{(\si{g/cm^2})}=\frac{X_{\max}}{X_0}\,X_0,
    \qquad X_0=\SI{37.8}{g/cm^2}.
\]
\end{calculationbox}

\begin{calculationbox}[title=Numerik]
F\"ur $10^{11}\,\si{eV}$:
\[
    X_{\max}=10.22\cdot \SI{37.8}{g/cm^2}=\SI{386.21}{g/cm^2},
\]
\[
    h=-\SI{8}{km}\ln\paren{\frac{386.21}{1030}}=\SI{7.85}{km}.
\]
F\"ur $10^{16}\,\si{eV}$:
\[
    X_{\max}=26.83\cdot \SI{37.8}{g/cm^2}=\SI{1014.06}{g/cm^2},
\]
\[
    h=-\SI{8}{km}\ln\paren{\frac{1014.06}{1030}}=\SI{0.12}{km}.
\]
\end{calculationbox}

\begin{answerbox}
\[
    \boxed{h(10^{11}\,\si{eV})\approx \SI{7.85}{km}},
    \qquad
    \boxed{h(10^{16}\,\si{eV})\approx \SI{0.12}{km}}.
\]
Der hochenergetischere Schauer entwickelt sich tiefer in der Atmosph\"are, weil mehr Verdopplungsschritte n\"otig sind, bis die Energie pro Teilchen auf $E_c$ gefallen ist.
\end{answerbox}

\subsection*{(e) Energie f\"ur ein Schauermaximum am Technik-Campus}

\begin{calculationbox}
Wir verwenden f\"ur den Technik-Campus n\"aherungsweise die H\"ohe von Innsbruck
\[
    h\approx \SI{0.574}{km}.
\]
Damit ist die atmosph\"arische Tiefe dort
\[
    X(h)=\SI{1030}{g/cm^2}\,e^{-0.574/8}=\SI{958.69}{g/cm^2}.
\]
In Einheiten von $X_0$ ergibt sich
\[
    \frac{X_{\max}}{X_0}=\frac{958.69}{37.8}=25.36.
\]
Aus
\[
    \frac{X_{\max}}{X_0}=\frac{\ln(E_0/E_c)}{\ln 2}
\]
folgt
\[
    E_0=E_c\,2^{X_{\max}/X_0}.
\]
Also
\[
    E_0=8.4\cdot 10^7\,\si{eV}\cdot 2^{25.36}
    =3.62\cdot 10^{15}\,\si{eV}.
\]
\end{calculationbox}

\begin{answerbox}
\[
    \boxed{E_0\approx 3.62\cdot 10^{15}\,\si{eV}=\SI{3.62}{PeV}}.
\]
Ein Elektron der kosmischen Strahlung m\"usste also eine Energie im PeV-Bereich haben, damit das Schauermaximum ungef\"ahr auf H\"ohe des Technik-Campus liegt.
\end{answerbox}

\begin{summarybox}[title=Tafelvortrag zu Aufgabe 1]
\begin{enumerate}
    \item Zuerst das Modell erkl\"aren: nach jeder Strahlungsl\"ange $X_0$ Verdopplung der Teilchenzahl und Halbierung der Energie pro Teilchen.
    \item Dann die zwei zentralen Formeln anschreiben:
    \[
        N(n)=2^n,\qquad E(n)=\frac{E_0}{2^n}.
    \]
    \item Das Maximum folgt aus $E(n_{\max})=E_c$.
    \item Die Umrechnung in eine H\"ohe erfolgt nicht direkt \"uber geometrische Strecke, sondern \"uber die atmosph\"arische Tiefe $X$ in $\si{g/cm^2}$.
\end{enumerate}
\end{summarybox}

\FloatBarrier
\newpage

%%%%%%%%%%%%%%%%%%%%%%%%%%%%%%%%%%%%%%%%%%%%%%%%%%%%%%%%%%%%%%%%%%%
\section{Aufgabe 2: Cherenkov-Licht in der Atmosph\"are}

\begin{taskbox}
Ein geladenes Teilchen emittiert Cherenkov-Licht, wenn es sich in einem Medium schneller bewegt als Licht in diesem Medium. Die Lichtgeschwindigkeit im Medium ist $c_n=c/n$. Der Emissionswinkel relativ zur Flugrichtung erf\"ullt
\[
    \cos\vartheta_c=\frac{1}{\beta n}.
\]
Gegeben ist
\[
    n(h)=1+0.000283\,e^{-h/H_0},
    \qquad H_0=\SI{8}{km}.
\]
\end{taskbox}

\begin{definitionbox}[title=Cherenkov-Bedingung]
Cherenkov-Licht entsteht nur, wenn
\[
    v>\frac{c}{n}
    \quad\Longleftrightarrow\quad
    \beta n>1.
\]
Ist $\beta n\leq 1$, ist der Ausdruck $1/(\beta n)$ gr\"o\ss{}er oder gleich $1$ und es gibt keinen realen Cherenkov-Winkel.
\end{definitionbox}

\subsection*{(a) Cherenkov-Winkel eines Positrons mit $E=\SI{200}{MeV}$ bei $h=\SI{10}{km}$}

\begin{calculationbox}
F\"ur ein Positron ist
\[
    m_ec^2=\SI{0.511}{MeV}.
\]
Wenn $E=\SI{200}{MeV}$ als Gesamtenergie verwendet wird, ist
\[
    \gamma=\frac{E}{m_ec^2}=\frac{200}{0.511}=391.39.
\]
Daraus folgt
\[
    \beta=\sqrt{1-\frac{1}{\gamma^2}}
    =\sqrt{1-\frac{1}{391.39^2}}
    =0.9999967.
\]
\end{calculationbox}

\begin{calculationbox}
Der Brechungsindex in $h=\SI{10}{km}$ ist
\[
    n(\SI{10}{km})=1+0.000283\,e^{-10/8}=1.0000811.
\]
Damit ist
\[
    \cos\vartheta_c=\frac{1}{\beta n}
    =\frac{1}{0.9999967\cdot 1.0000811}
    =0.9999222.
\]
Also
\[
    \vartheta_c=\arccos(0.9999222)
    =0.01247\,\si{rad}
    =0.714^\circ.
\]
\end{calculationbox}

\begin{answerbox}
\[
    \boxed{\vartheta_c\approx 0.01247\,\si{rad}\approx 0.714^\circ}.
\]
\end{answerbox}

\begin{notebox}
W\"urde $E=\SI{200}{MeV}$ stattdessen als kinetische Energie interpretiert, w\"are $\gamma=(T+m_ec^2)/(m_ec^2)$ minimal gr\"o\ss{}er. Der Winkel \"andert sich hier praktisch nicht, weil das Positron ohnehin ultrarelativistisch ist.
\end{notebox}

\subsection*{(b) Radius des Cherenkov-Kegels auf Meeresh\"ohe}

\begin{calculationbox}
Das Positron komme direkt aus dem Zenit, also senkrecht nach unten. Dann ergibt die einfache Geometrie eines Kegels
\[
    r=h\tan\vartheta_c.
\]
Mit
\[
    h=\SI{10}{km}=\SI{10000}{m},
    \qquad
    \vartheta_c=0.01247\,\si{rad}
\]
folgt
\[
    r=\SI{10000}{m}\tan(0.01247)=\SI{124.7}{m}.
\]
\end{calculationbox}

\begin{answerbox}
\[
    \boxed{r\approx \SI{125}{m}}.
\]
\end{answerbox}

\begin{summarybox}[title=Tafelvortrag zu Aufgabe 2]
\begin{enumerate}
    \item Erst die Schwelle $\beta n>1$ nennen.
    \item Dann aus der Positronenergie $\gamma$ und $\beta$ berechnen.
    \item Danach den Brechungsindex in $\SI{10}{km}$ einsetzen.
    \item Der Radius folgt geometrisch aus $r=h\tan\vartheta_c$.
\end{enumerate}
\end{summarybox}

\FloatBarrier
\newpage

%%%%%%%%%%%%%%%%%%%%%%%%%%%%%%%%%%%%%%%%%%%%%%%%%%%%%%%%%%%%%%%%%%%
\section{Aufgabe 3: Energieverlust durch Ionisation und Bethe-Bloch-Formel}

\begin{taskbox}
F\"ur schwere geladene Teilchen beschreibt die Bethe-Bloch-Formel den mittleren Energieverlust durch Ionisation und Anregung des Mediums:
\[
-\dv{E}{x}
=\frac{1}{(4\pi\epsilon_0)^2}\frac{4\pi e^4z^2}{m_ec^2}N_A\frac{Z}{A}\rho\frac{1}{\beta^2}
\paren{
\ln\paren{\frac{2m_ec^2\beta^2\gamma^2}{I}}-\beta^2
}.
\]
\end{taskbox}

\begin{definitionbox}[title=Bedeutung der wichtigsten Faktoren]
\begin{itemize}
    \item $z$: Ladungszahl des Projektils. Der Energieverlust skaliert stark mit $z^2$.
    \item $\beta=v/c$: langsame Teilchen ionisieren stark, weil der Faktor $1/\beta^2$ gro\ss{} wird.
    \item $Z,A,\rho,I$: Materialparameter des Absorbers.
    \item Der Logarithmus erzeugt bei hohen Energien den langsamen relativistischen Anstieg.
\end{itemize}
\end{definitionbox}

\subsection*{(a) Energieverlust in einem \SI{2}{cm} dicken Graphitabsorber}

\begin{calculationbox}[title=Eingangsgr\"o\ss{}en]
Gegeben bzw. verwendet:
\[
\begin{array}{lll}
T=\SI{1000}{MeV}, & d=\SI{2}{cm}, & m_ec^2=\SI{0.511}{MeV},\\
 m_pc^2=\SI{938.3}{MeV}, & m_Dc^2=\SI{1875.6}{MeV}, & m_\alpha c^2=\SI{3727.4}{MeV},\\
 Z_C=6, & A_C=\SI{12.01}{g/mol}, & \rho=\SI{2.2}{g/cm^3}.
\end{array}
\]
F\"ur Graphit wird das mittlere Anregungspotential gesch\"atzt durch
\[
    I\approx 16\,Z^{0.9}\,\si{eV}=16\cdot 6^{0.9}\,\si{eV}=\SI{80.25}{eV}
    =8.025\cdot 10^{-5}\,\si{MeV}.
\]
Praktisch verwendet man
\[
    K=\frac{1}{(4\pi\epsilon_0)^2}\frac{4\pi e^4}{m_ec^2}N_A
    \approx \SI{0.307075}{MeV.cm^2/mol}.
\]
Dann wird
\[
    -\dv{E}{x}
    =Kz^2\rho\frac{Z}{A}\frac{1}{\beta^2}
    \paren{\ln\paren{\frac{2m_ec^2\beta^2\gamma^2}{I}}-\beta^2}.
\]
\end{calculationbox}

\begin{calculationbox}[title=Relativistische Kinematik]
F\"ur jedes Teilchen mit kinetischer Energie $T$ gilt
\[
    \gamma=\frac{T+mc^2}{mc^2}=1+\frac{T}{mc^2},
    \qquad
    \beta=\sqrt{1-\frac{1}{\gamma^2}}.
\]
Die berechneten Werte sind:
\[
\begin{array}{c|c|c|c|c|c}
\text{Teilchen} & z & \gamma & \beta & -\dd E/\dd x\,[\si{MeV/cm}] & \Delta E\,[\si{MeV}]\\
\hline
p & 1 & 2.0658 & 0.8750 & 4.35 & 8.70\\
D & 1 & 1.5332 & 0.7580 & 5.39 & 10.78\\
\alpha & 2 & 1.2683 & 0.6151 & 30.61 & 61.21
\end{array}
\]
Dabei ist
\[
    \Delta E=\paren{-\dv{E}{x}}d,
    \qquad d=\SI{2}{cm}.
\]
\end{calculationbox}

\begin{answerbox}
\[
    \boxed{\Delta E_p\approx\SI{8.70}{MeV}},
    \qquad
    \boxed{\Delta E_D\approx\SI{10.78}{MeV}},
    \qquad
    \boxed{\Delta E_\alpha\approx\SI{61.21}{MeV}}.
\]
Das $\alpha$-Teilchen verliert deutlich mehr Energie, weil es die Ladung $z=2$ besitzt und der Energieverlust daher einen Faktor $z^2=4$ enth\"alt. Zus\"atzlich ist es bei gleicher kinetischer Energie langsamer, wodurch $1/\beta^2$ gr\"o\ss{}er wird.
\end{answerbox}

\subsection*{(b) Qualitative Skizze von $-\dd E/\dd x$ von \SI{1}{MeV} bis \SI{1}{TeV}}

\begin{calculationbox}[title=Form der Bethe-Bloch-Kurve]
Die Kurve hat drei charakteristische Bereiche:
\begin{enumerate}
    \item \textbf{Niedrige Energien:} $\beta$ ist klein und $1/\beta^2$ ist gro\ss{}. Der Energieverlust ist stark.
    \item \textbf{Minimum:} Bei mittleren Energien erreicht $-\dd E/\dd x$ ein Minimum. Teilchen in diesem Bereich hei\ss{}en minimal ionisierende Teilchen.
    \item \textbf{Hohe Energien:} Der Logarithmus erzeugt einen langsamen relativistischen Anstieg.
\end{enumerate}
\end{calculationbox}

\begin{calculationbox}[title=Qualitative log-log-Skizze]
\begin{center}
\begin{tikzpicture}
\begin{loglogaxis}[
    width=0.86\textwidth,
    height=7cm,
    xlabel={$T\,[\si{MeV}]$},
    ylabel={$-\dd E/\dd x\,[\si{MeV/cm}]$},
    xmin=1, xmax=1e6,
    ymin=1, ymax=100,
    grid=both,
    legend pos=south west,
]
\addplot[thick,domain=1:1e6,samples=250] {2.6 + 35/(x^0.75+1) + 0.20*ln(x+1)};
\addlegendentry{qualitativ}
\addplot[only marks,mark=*] coordinates {(1000,4.35)};
\addlegendentry{$p$ bei \SI{1}{GeV}}
\addplot[only marks,mark=square*] coordinates {(1000,5.39)};
\addlegendentry{$D$ bei \SI{1}{GeV}}
\addplot[only marks,mark=triangle*] coordinates {(1000,30.61)};
\addlegendentry{$\alpha$ bei \SI{1}{GeV}}
\end{loglogaxis}
\end{tikzpicture}
\end{center}
\end{calculationbox}

\begin{answerbox}
Die qualitative Kurve f\"allt bei kleinen Energien stark ab, erreicht ein Minimum und steigt bei sehr hohen Energien nur langsam wieder an. Die Werte aus Teil (a) dienen als Normierungspunkte bei $T=\SI{1}{GeV}$.
\end{answerbox}

\subsection*{(c) Gilt die Bethe-Bloch-Formel auch f\"ur Elektronen?}

\begin{calculationbox}
Die angegebene Bethe-Bloch-Formel ist in dieser Form f\"ur schwere geladene Teilchen wie Myonen, Protonen, Deuteronen oder $\alpha$-Teilchen gedacht. F\"ur Elektronen ist die Situation anders, weil das Projektil dieselbe Masse wie die Target-Elektronen hat.

Daraus folgen zwei wichtige \"Anderungen:
\begin{itemize}
    \item Die Kinematik der Energie\"ubertragung ist anders als bei schweren Projektilen.
    \item Projektil-Elektronen und Target-Elektronen sind ununterscheidbare Teilchen; das muss in der Streuformel ber\"ucksichtigt werden.
\end{itemize}
Zus\"atzlich werden bei Elektronen Strahlungsverluste durch Bremsstrahlung schon deutlich fr\"uher wichtig als bei schweren Teilchen.
\end{calculationbox}

\begin{answerbox}
\[
    \boxed{\text{Nicht direkt in dieser einfachen Form.}}
\]
F\"ur Elektronen braucht man eine modifizierte Behandlung der Kollisionsverluste und zus\"atzlich die Ber\"ucksichtigung von Bremsstrahlung. Die Bethe-Bloch-Idee -- Energieverlust durch Ionisation und Anregung -- bleibt relevant, aber die konkrete Formel ist nicht unver\"andert anwendbar.
\end{answerbox}

\begin{summarybox}[title=Tafelvortrag zu Aufgabe 3]
\begin{enumerate}
    \item Zuerst klarmachen: Bethe-Bloch beschreibt mittleren Ionisationsverlust schwerer geladener Teilchen.
    \item Dann alle Eingangsgr\"o\ss{}en nennen: $T,m,z,Z,A,\rho,I,d$.
    \item Wichtigster Vergleich: $\alpha$ verliert viel mehr Energie wegen $z^2=4$ und kleinerem $\beta$.
    \item Bei der Kurve die drei Zonen erkl\"aren: $1/\beta^2$-Abfall, Minimum, logarithmischer Anstieg.
    \item F\"ur Elektronen: nicht direkt dieselbe Formel, weil gleiche Masse wie Target-Elektronen und Bremsstrahlung wichtiger.
\end{enumerate}
\end{summarybox}

\FloatBarrier
\newpage

%%%%%%%%%%%%%%%%%%%%%%%%%%%%%%%%%%%%%%%%%%%%%%%%%%%%%%%%%%%%%%%%%%%
\section{Lernzusammenfassung f\"ur das gesamte Blatt}

\begin{summarybox}[title=Die wichtigsten Formeln]
\begin{align*}
    N(n)&=2^n, & E(n)&=\frac{E_0}{2^n},\\
    \frac{X_{\max}}{X_0}&=\frac{\ln(E_0/E_c)}{\ln 2}, & N_{\max}&=\frac{E_0}{E_c},\\
    X(h)&=X_Ee^{-h/H_0}, & h&=-H_0\ln\paren{\frac{X}{X_E}},\\
    \cos\vartheta_c&=\frac{1}{\beta n}, & r&=h\tan\vartheta_c,\\
    \gamma&=1+\frac{T}{mc^2}, & \beta&=\sqrt{1-\frac{1}{\gamma^2}}.
\end{align*}
\end{summarybox}

\begin{warningbox}[title=Typische Fehlerquellen]
\begin{itemize}
    \item $X$ in $\si{g/cm^2}$ ist keine geometrische L\"ange, sondern eine atmosph\"arische Tiefe.
    \item Beim Cherenkov-Effekt muss zuerst $\beta n>1$ erf\"ullt sein.
    \item Bei Bethe-Bloch nicht vergessen: $I$ muss in derselben Energieeinheit wie $m_ec^2$ eingesetzt werden.
    \item $\Delta E$ durch einen Absorber ist $(-\dd E/\dd x)d$; die Formel liefert zuerst einen Energieverlust pro L\"ange.
    \item Die einfache Bethe-Bloch-Formel gilt nicht unver\"andert f\"ur Elektronen.
\end{itemize}
\end{warningbox}

\begin{notebox}[title=Mini-Check vor einem Tafelvortrag]
Man sollte die folgenden S\"atze frei sagen k\"onnen:
\begin{enumerate}
    \item Ein elektromagnetischer Schauer ist ein Wechselspiel aus Bremsstrahlung und Paarbildung.
    \item Das Schauermaximum ist erreicht, wenn die mittlere Teilchenenergie auf $E_c$ gefallen ist.
    \item Cherenkov-Licht entsteht, wenn ein geladenes Teilchen schneller ist als Licht im Medium.
    \item Bethe-Bloch erkl\"art den Ionisationsverlust schwerer geladener Teilchen und zeigt den typischen Verlauf: hoher Verlust bei kleinen Energien, Minimum, relativistischer Anstieg.
\end{enumerate}
\end{notebox}

\end{document}
